\documentclass[12pt,a4paper,oneside]{report}

% --- PACKAGES ---
\usepackage[utf8]{inputenc}
\usepackage[italian]{babel}
\usepackage[T1]{fontenc}
\usepackage{graphicx} % Per le immagini
\usepackage{geometry} % Per i margini
 \geometry{margin=2.5cm}
\usepackage{hyperref} % Per i link cliccabili
\usepackage{amsmath, amssymb} % Per le formule matematiche
\usepackage{listings} % Per il codice Processing
\usepackage{xcolor} % Per i colori del codice
\usepackage{fancyhdr} % Per testata e piè di pagina
\usepackage{longtable, booktabs}

% --- CONFIGURAZIONE CODICE (Processing/Java) ---
\lstset{
    language=Java,
    basicstyle=\ttfamily\small,
    keywordstyle=\color{blue},
    stringstyle=\color{red},
    commentstyle=\color{green!50!black},
    breaklines=true,
    frame=single
}

% --- COMANDI EXTRA PER PANDOC ---
\providecommand{\tightlist}{%
  \setlength{\itemsep}{0pt}\setlength{\parskip}{0pt}}

\usepackage{calc} 
\usepackage{longtable,booktabs,array}

\makeatletter
\newsavebox\pandoc@box
\newcommand*\pandocbounded[1]{% scales image to fit in text height/width
  \sbox\pandoc@box{#1}%
  \Gscale@div\@tempa{\textheight}{\dimexpr\ht\pandoc@box+\dp\pandoc@box\relax}%
  \Gscale@div\@tempb{\linewidth}{\wd\pandoc@box}%
  \ifdim\@tempb\p@<\@tempa\p@\let\@tempa\@tempb\fi%
  \ifdim\@tempa\p@<\p@\scalebox{\@tempa}{\usebox\pandoc@box}%
  \else\usebox{\pandoc@box}%
  \fi%
}
\makeatother

\begin{document}

% 1. Copertina
\input{copertina/copertina.tex}

\pagenumbering{roman} % Numerazione romana per premessa e indice

% 2. Premessa
\input{premessa/premessa.tex}

% 3. Info Esame, per adesso non esiste
%\input{premessa/info_esame.tex}

% 4. Indice
\tableofcontents
\newpage

\pagenumbering{arabic} % Torna alla numerazione normale

% 5. Teoria (Lezioni 1-12)
\part{Teoria}
Osservatori Standard e le CMF

Poiché ogni persona vede i colori in modo leggermente diverso a causa della sensibilità della propria retina, la scienza ha creato un Osservatore Standard, ovvero un "occhio matematico" medio.

Le CMF (Color-Matching Functions) sono tre curve, chiamate $\overline{x}(\lambda)$, $\overline{y}(\lambda)$ e $\overline{z}(\lambda)$, che descrivono la sensibilità spettrale di questo osservatore ideale.

- Queste curve rappresentano come l'occhio risponde alle componenti del rosso, del verde e del blu per ogni lunghezza d'onda.
- Vengono utilizzate per calcolare i valori di tristimolo X, Y, Z, che sono le coordinate numeriche finali di un colore.

### CIE 1931 e CIE 1964: Qual è la differenza?

La differenza principale tra questi due standard riguarda quanto spazio occupa l'oggetto che stiamo guardando nel nostro campo visivo, espresso come angolo di osservazione.

- CIE 1931 (Osservatore a 2°): È stato progettato studiando la visione umana su un angolo di soli 2 gradi. Immagina di guardare la punta di un dito a circa 50 cm di distanza: in quell'area piccolissima della retina (la fovea) siamo molto sensibili ai dettagli, ma vediamo meno superficie. È adatto per oggetti visti da lontano o molto piccoli.
- CIE 1964 (Osservatore supplementare a 10°): Poiché l'occhio si comporta diversamente quando guarda superfici più ampie, nel 1964 è stato introdotto un osservatore basato su un angolo di 10 gradi. Immagina di guardare un oggetto grande quanto un'arancia a 50 cm di distanza. Questo standard è più accurato per misurare il colore di superfici estese. % Assicurati che il percorso sia identico a quello del file tree
\input{Teoria/Lezione 2/lezione2.tex}
\input{Teoria/Lezione 3/lezione3.tex}
\input{Teoria/Lezione 4/lezione4.tex}
\input{Teoria/Lezione 5/lezione5.tex}
\input{Teoria/Lezione 6/lezione6.tex}
\input{Teoria/Lezione 7/lezione7.tex}
\input{Teoria/Lezione 8/lezione8.tex}
\input{Teoria/Lezione 9/lezione9.tex}
\input{Teoria/Lezione 10/lezione10.tex}
\input{Teoria/Lezione 11/lezione11.tex}
\input{Teoria/Lezione 12/lezione12.tex}
% Aggiungi qui gli altri \input man mano che crei i file

% 6. Laboratorio (Processing)
\part{Laboratorio}
\input{laboratorio/lezione1_lab.tex}
% Aggiungi qui gli altri \input

% 7. Formule
\appendix
\input{formule/formule.tex}

\end{document}